\documentclass[11pt]{article}

\usepackage[letterpaper,margin=1in]{geometry}
\usepackage[T1]{fontenc}
\usepackage[utf8]{inputenc}
\usepackage{lmodern}
\usepackage{microtype}
\usepackage{amsmath,amssymb,amsthm,mathtools}
\usepackage{booktabs,longtable,tabularx,array}
\usepackage{enumitem}
\usepackage{xcolor}
\usepackage{graphicx}
\usepackage{tikz}
\usetikzlibrary{arrows.meta,positioning,calc}
\usepackage{fancyhdr}
\usepackage{lastpage}
\usepackage{hyperref}
\usepackage[nameinlink,capitalise,noabbrev]{cleveref}

\definecolor{deepblue}{RGB}{19,63,104}
\definecolor{teal}{RGB}{0,108,112}
\definecolor{brick}{RGB}{146,54,44}
\definecolor{softgray}{RGB}{245,247,249}
\definecolor{midgray}{RGB}{92,100,110}

\hypersetup{
  colorlinks=true,
  linkcolor=deepblue,
  citecolor=teal,
  urlcolor=brick,
  pdftitle={From Finite Blacklists to Freiling Symmetry},
  pdfauthor={Research report prepared for Vladimir Reshetnikov},
  pdfsubject={Finite and infinite sparse symmetry avoidance},
  pdfkeywords={axiom of choice, Freiling symmetry, graph coloring, independent sets, tournaments, descriptive set theory}
}

\pagestyle{fancy}
\fancyhf{}
\fancyhead[L]{\small\textsc{Sparse Symmetry Avoidance}}
\fancyhead[R]{\small\textsc{Finite, Infinite, and Definable Forms}}
\fancyfoot[C]{\small Page \thepage\ of \pageref*{LastPage}}
\setlength{\headheight}{14pt}

\newtheorem{theorem}{Theorem}[section]
\newtheorem{proposition}[theorem]{Proposition}
\newtheorem{lemma}[theorem]{Lemma}
\newtheorem{corollary}[theorem]{Corollary}
\newtheorem{conjecture}[theorem]{Conjecture}
\theoremstyle{definition}
\newtheorem{definition}[theorem]{Definition}
\newtheorem{problem}[theorem]{Problem}
\newtheorem{question}[theorem]{Research question}
\theoremstyle{remark}
\newtheorem{remark}[theorem]{Remark}
\newtheorem{example}[theorem]{Example}

\newcommand{\cF}{\mathcal F}
\newcommand{\cA}{\mathcal A}
\newcommand{\cC}{\mathcal C}
\newcommand{\cP}{\mathcal P}
\newcommand{\N}{\mathbb N}
\newcommand{\R}{\mathbb R}
\newcommand{\Z}{\mathbb Z}
\newcommand{\ind}{\mathbf 1}
\newcommand{\Safe}{\operatorname{Safe}}
\newcommand{\Conf}{\operatorname{Conf}}
\newcommand{\BPI}{\mathsf{BPI}}
\newcommand{\SAP}{\mathsf{SAP}}
\newcommand{\Cm}{\mathsf{C}}
\newcommand{\vcent}{\mathbin{:}}
\newcommand{\defeq}{\vcent=}
\newcommand{\card}[1]{\lvert #1\rvert}
\newcommand{\set}[1]{\left\{#1\right\}}
\newcommand{\angles}[1]{\left\langle#1\right\rangle}
\newcommand{\eps}{\varepsilon}
\newcommand{\doi}[1]{\href{https://doi.org/#1}{doi:#1}}
% Added by the ProveIt write phase (batch 114, 5 October 2026): dated notes
% marked [write].
\newcommand{\file}[1]{\nolinkurl{#1}}
\newenvironment{writenote}[1][]{\par\smallskip\begingroup\small
 \noindent\textbf{[write]}\if\relax\detokenize{#1}\relax\else~\textbf{#1.}\fi\ }%
 {\par\endgroup\smallskip}

\setlist[itemize]{leftmargin=1.5em,itemsep=0.35em,topsep=0.35em}
\setlist[enumerate]{leftmargin=1.7em,itemsep=0.35em,topsep=0.35em}
\renewcommand{\arraystretch}{1.16}

\begin{document}

\hypersetup{pageanchor=false}
\begin{titlepage}
\centering
\vspace*{0.65in}
{\color{deepblue}\rule{0.88\textwidth}{1.4pt}}\\[0.35in]
{\Huge\bfseries From Finite Blacklists to Freiling Symmetry\par}
\vspace{0.15in}
{\LARGE Exact avoidance thresholds, critical tournaments,\\
choice-sensitive colorings, and perfect safe sets\par}
\vspace{0.33in}
{\large A finite--infinite research problem at the intersection of\\
extremal combinatorics, weak choice, measure, and topology\par}
\vspace{0.45in}
{\color{deepblue}\rule{0.88\textwidth}{1.4pt}}\\[0.55in]
{\Large Research report prepared for Vladimir Reshetnikov\par}
\vspace{0.12in}
{\large 5 October 2026\par}
\vfill
\begin{minipage}{0.88\textwidth}
\small
\noindent\textbf{Status and attribution.}
This is an AI-assisted, unrefereed research report.  It proposes a problem,
proves the results stated as theorems below, and separates those results from
classical ingredients and conjectural extensions.  The elementary infinite
partition theorem is due to de Bruijn and Erd\H{o}s (1951), not to this report.
No exhaustive priority search has been made for the finite refinements proved
here.  No endorsement, collaboration, or authorship by Elliot Glazer is
implied; the motivation is based only on publicly available work.
\end{minipage}
\end{titlepage}
\hypersetup{pageanchor=true}
\pagenumbering{arabic}

\begin{abstract}
A \emph{blacklist system} assigns to each point $x$ a set $F(x)$ of points
that $x$ refuses to pair with.  A set is \emph{safe} when no two of its
members blacklist one another.  This simple relation problem is a finite,
quantitative cousin of Freiling's axiom of symmetry and a natural bridge
between extremal graph theory and choiceless mathematics.

For $n$ points and blacklists of size at most $k$, put $q=2k+1$.  We prove
that every system has a safe set of size $\lceil n/q\rceil$, and that this is
best possible for every $n$.  Equivalently, a safe $r$-set is forced at the
exact threshold
\[
     n=(r-1)(2k+1)+1.
\]
At the last failing value $n=(r-1)q$, every extremal system is rigid: its
conflict graph is a disjoint union of $r-1$ copies of $K_q$, and each block is
oriented as a regular tournament.  This gives an exact enumeration of all
labelled critical systems.  We also determine the minimum number of safe
pairs exactly, obtain the first two asymptotic orders for the number of safe
$r$-sets for every fixed $k,r$, prove an exact formula for safe triples when
$k=1$, and establish a heterogeneous list-colouring theorem with local list
size $2\lvert F(x)\rvert+1$.

The infinite form exposes a foundational boundary.  The classical
compactness theorem of de Bruijn--Erd\H{o}s yields a partition into $q$ safe
sets under a suitable choice principle.  Over ZF we isolate the principle
$\SAP_k$ and prove
\[
   \BPI\Longrightarrow\SAP_k
   \Longrightarrow \bigwedge_{2\le m\le k+1}\Cm_m,
\]
where $\BPI$ is the Boolean Prime Ideal Theorem and $\Cm_m$ is choice for
families of $m$-element sets.  We also prove $\SAP_k$ for every well-orderable
underlying set in ZF.  In the definable setting, regularity gives far more:
a Borel relation with countable vertical sections on a perfect Polish space
has a perfect safe set, while a measurable such relation on an atomless
standard probability space admits an almost surely safe countably infinite
i.i.d. sample.  These results turn the proposed problem into a concrete
research program about exact finite spectra, weak choice, Borel colourings,
and generalized symmetry axioms.
\end{abstract}

\begin{writenote}[Status and trust boundary, added 5 October 2026, batch~114]
This is an AI-assisted, unrefereed external manuscript (its status box says
so), filed in ProveIt as the single-source research report
\file{freiling-symmetry-blacklists} (provenance, the sources as the write
read them, relation to the repository, the collected non-claims and the
reading conventions are in Section~\ref{fbl:sec:provenance}). Its author
line reads ``Research report prepared for Vladimir Reshetnikov''. No
statement of the report is formalized in Lean or Rocq; the formalization
plan of Section~\ref{fbl:sec:formal} is a proposal.
\emph{What rests on what.} The partition of an arbitrary infinite system
into $2k+1$ safe sets is Theorem~3 of de Bruijn and
Erd\H{o}s~\cite{deBruijnErdos}, which the write read: the finite case of
Corollary~\ref{fbl:cor:qcolor} and its low-degree-vertex induction are
already there (note after Corollary~\ref{fbl:cor:weighted}). Tur\'an's
theorem with its equality case is quoted without proof in
Theorem~\ref{fbl:thm:critical}; propositional compactness from $\BPI$ in
Theorem~\ref{fbl:thm:dbe}; the Kuratowski--Ulam and Mycielski theorems in
Theorem~\ref{fbl:thm:perfect}; Tonelli's theorem in
Theorem~\ref{fbl:thm:random}. Every other theorem, corollary, lemma and
proposition is proved in the text; the intake and the write reread every
proof and found no error. The exhaustive program \file{code/verification.py}
checks the case $k=1$ for $n\le7$, the count $40$ by enumeration and the
count $72{,}576$ only by evaluating \eqref{fbl:eq:critical-count} (note in
Section~\ref{fbl:sec:computation}); it proves nothing beyond those cases. The
source makes no historical-priority claim for its finite refinements
``without a more exhaustive literature review'' (delivery README), and the
write made no such review either. The write adds the notes marked [write],
Section~\ref{fbl:sec:provenance} and Section~\ref{fbl:sec:further} with
three propositions proved there.
\end{writenote}

\tableofcontents
\newpage

\section{Why this problem}

\subsection{The intersection of interests}
Elliot Glazer's dissertation studies interactions among fragments of the
Axiom of Choice, forcing, measure theory, box-game paradoxes, and weak
symmetry principles, explicitly including Freiling's axiom of symmetry
\cite{GlazerThesis}.  His paper on a choiceless box-game paradox asks how much
of the familiar predictive phenomenon survives when choice is removed
\cite{GlazerBox}.  His topological Tennenbaum theorem is a separate example
of extracting strong mathematical consequences from topological and Borel
regularity assumptions \cite{GlazerTennenbaum}.

The \texttt{ProveIt} repository already contains substantial work on
measurable box games and infinite hat guessing \cite{ProveIt,ProveItBoxes}.
A useful new direction should therefore remain close to choice, symmetry,
and definability without merely repeating that material.  Sparse
symmetry-avoidance does exactly this:

\begin{center}
\begin{tabularx}{0.93\textwidth}{>{\bfseries}p{0.23\textwidth}X}
Finite face & exact extremal thresholds, equality cases, enumeration,
independence polynomials, and algorithms;\\
Infinite face & compactness, weak choice, and the strength of a global
finite partition;\\
Definable face & Borel category, atomless measure, perfect sets, and a
regularity-based strengthening of Freiling's conclusion.\\
\end{tabularx}
\end{center}

\subsection{Freiling symmetry in one line}\label{fbl:sec:freiling}
Freiling's axiom of symmetry concerns maps
$f:\R\to[\R]^{\le\aleph_0}$ and asserts the existence of $x,y$ such that
$y\notin f(x)$ and $x\notin f(y)$ \cite{Freiling}.  In the language of this
report, it asks for a safe pair when every point has a countable blacklist.
The finite problem below replaces ``countable'' by a fixed bound $k$, asks
for much larger safe sets and exact counts, and then returns to infinite sets
to study global partitions and definable relations.

\begin{writenote}[Freiling and Sierpi\'nski, 5 October 2026]
The write confirmed the data of~\cite{Freiling} (J. Symbolic Logic
\textbf{51}, no.~1, March 1986, 190--200, doi:10.2307/2273955) and read its
abstract and reference list on the publisher's page, not its body. The
abstract presents the axioms as based on ``some intuition of Stuart Davidson
and an old theorem of Sierpi\'nski''; the reference list cites
Sierpi\'nski's book~\cite{Sierpinski}, which the write did not read. The
statement of the axiom above and its well-known equivalence with the
negation of the continuum hypothesis in ZFC, usually credited to
Sierpi\'nski, are therefore not checked here against Freiling's text or
Sierpi\'nski's. This report proves only one direction of that
equivalence, and only in the form of the remark after
Corollary~\ref{fbl:cor:freiling}: under CH a well-ordering of $\R$ in type
$\omega_1$ gives countable blacklists with no safe pair (each of two
distinct reals precedes the other or follows it). The converse, that
Freiling's axiom holds when CH fails, is neither stated nor used in this
report.
\end{writenote}

\subsection{The classical theorem that must be credited}\label{fbl:sec:dbe-credit}
De Bruijn and Erd\H{o}s proved in 1951 that if $\lvert F(x)\rvert\le k$ for
all $x$, then the entire set is the union of $2k+1$ safe sets
\cite{deBruijnErdos}.  Their paper explicitly adopts the Axiom of Choice and
derives the infinite statement from graph-colouring compactness.  The basic
infinite partition is therefore classical.  The point of this report is to
refine the finite theory sharply, expose the exact extremizers, quantify the
safe-set spectrum, and isolate the weak-choice and regularity questions that
the classical formulation leaves hidden.

\begin{writenote}[de Bruijn and Erd\H{o}s, read 5 October 2026]
The write read the paper~\cite{deBruijnErdos} in full (the scan of the
reprint in the Eindhoven University of Technology repository). It confirms
the paragraph above: ``The Axiom of Choice is adopted throughout the
paper'', and its Theorem~1 (a graph all of whose finite subgraphs are
$k$-colourable is $k$-colourable) is deduced from Rado's selection theorem,
with a proof by Tychonoff's theorem credited to Rabson and A.~Stone. Their
Theorem~3 is Theorem~\ref{fbl:thm:dbe} below: for $f(b)\subseteq S\setminus
\{b\}$ with $\lvert f(b)\rvert\le k$, $S$ is the union of $2k+1$
\emph{independent} sets (their word for safe). Its proof for finite $S$ is
the argument of Lemma~\ref{fbl:lem:edge-budget} and
Corollary~\ref{fbl:cor:qcolor}: the graph has at most $k\lvert S\rvert$
edges, hence a vertex joined to fewer than $2k+1$ others, and induction on
$\lvert S\rvert$; the infinite case follows from Theorem~1. Their Theorem~4
is the countable-colour statement of Question~\ref{fbl:q:countable}: if every
$f(b)$ is finite, $S$ is the union of countably many independent sets
(apply Theorem~3 to each $S_k=\{b:\lvert f(b)\rvert=k\}$). The paper says that
Theorems~1, 3 and~4 were announced in Erd\H{o}s~\cite{Erdos1950}, which the
write did not read. In the language of combinatorial set theory a map $F$
with $x\notin F(x)$ is a \emph{set mapping} and a safe set a \emph{free
set}; Hajnal's free set theorem~\cite{Hajnal} (not read) belongs to that
literature. The source does not use this vocabulary, and the write did not
survey that literature; this is one more reason to keep the source's
priority caveat as it stands.
\end{writenote}

\subsection{Provenance, sources and reading conventions}
\label{fbl:sec:provenance}
\begin{writenote}[Added 5 October 2026, batch~114]
\emph{Provenance.} The report is built from one manuscript, the archive
\file{From_Finite_Blacklists_to_Freiling_Symmetry.zip} (wrapper directory of
the same name, 6 files), which reached ProveIt's \file{docs/incoming} in the
arrival commit \file{2399df2bd} (5~October 2026, ``New research reports''),
outside the session bundle of the intake, and was placed as batch~114 in
\file{99053b5d1}. The text is the delivered
\file{From_Finite_Blacklists_to_Freiling_Symmetry.tex} (1{,}291 lines),
shipped as \file{article.tex}; the delivered \file{verification.py} and
\file{verification_output.txt} are shipped unchanged in \file{code/} and
\file{data/}. The delivered 19-page PDF and \file{SHA256SUMS.txt} (four
entries, all verified at placement) are not shipped; they and the delivery
\file{README.txt} (staged at placement and replaced by the report README)
are retrievable from the arrival commit. The package pins ProveIt revision
\file{a763feee1f3703dc710d9b810b1de66168f1e266} (bibliography
entry~\cite{ProveIt}); that commit is the batch-98A write commit of this
repository and an ancestor of the commit \file{9ccf04eae} this write is
based on. The status box calls the report ``AI-assisted'' and
``unrefereed''; the author line reads ``Research report prepared for
Vladimir Reshetnikov''. The package contains no e-mail address, no
``prepared for private review'' line and no personal data. Single source:
no merge choice was made. The write prefixed the 36 delivered labels with
\file{fbl:} and updated every reference; it added the notes marked [write],
this subsection, labels on eight delivered headings, one corollary, one
conjecture and the thirteen questions (none had a label), four
bibliography entries after the delivered ones, and
Section~\ref{fbl:sec:further} with
Propositions~\ref{fbl:prop:small}--\ref{fbl:prop:finite}. No statement,
proof or number of the source was changed: the added propositions are the
last numbered statements of their section and the added displays are
unnumbered, so every delivered theorem and equation keeps its number.

\emph{The sources, as the write read them} (5~October 2026).
\begin{itemize}\raggedright
\item de Bruijn and Erd\H{o}s~\cite{deBruijnErdos}: read in full (note
 above). The reprint is headed ``Reprinted from Proceedings, Series A, 54,
 No.~5 and Indag. Math., 13, No.~5, 1951'' with the single pagination
 371--373, and the paper was communicated at the meeting of 24~November
 1951. The bibliography's ``Indagationes Mathematicae \textbf{13} (1951),
 369--373'' is not supported by that copy, which gives 371--373 for both.
\item Freiling~\cite{Freiling}: data, abstract and reference list read, body
 not read (note in Section~\ref{fbl:sec:freiling}).
\item L\"auchli~\cite{Lauchli} and Mycielski~\cite{Mycielski}: bibliographic
 data confirmed through Crossref; the papers were not read. L\"auchli's
 paper is in the delivered bibliography but cited nowhere in the text; see
 Question~\ref{fbl:q:hierarchy} in Section~\ref{fbl:sec:further}.
\item Not read: Glazer's dissertation and papers~\cite{GlazerThesis,
 GlazerBox,GlazerTennenbaum} (motivation only), Cowen~\cite{Cowen1990,
 Cowen1998}, Howard and Rubin~\cite{HowardRubin}, Kechris~\cite{Kechris},
 and the uncited entries Tur\'an~\cite{Turan}, Jech~\cite{Jech} and
 Bollob\'as~\cite{Bollobas}. The standard facts the proofs take from them
 (propositional compactness from $\BPI$, the Kuratowski--Ulam and Mycielski
 theorems, Tur\'an's theorem and its equality case) are used as the source
 states them. [5~October 2026: Theorem~7.16 of Jech's book and the notes to
 its Chapter~7 were read when the independent check of the write was
 recorded; Jech is now cited in item~\ref{fbl:fq:sap1} of
 Section~\ref{fbl:sec:further}.]
\item The OEIS: the counts of distinct conflict graphs in
 Section~\ref{fbl:sec:computation} are terms of A133686, labelled graphs
 whose components have at most one cycle~\cite{OEIS}, and
 $R_1,R_3,R_5=1,2,24$ are terms of A007079, labelled regular tournaments
 on $2k+1$ nodes. The write identified both by searching the OEIS; the
 source names neither. OEIS data are licensed CC BY-SA~4.0; nothing was
 submitted.
\end{itemize}

\emph{What was checked at intake.} The batch-114 dossier read the
manuscript in full and checked the block construction, the Tur\'an count
$\binom t2q^2$, the arc-saturation step, formula~\eqref{fbl:eq:i3-formula},
the component costs in the proof of Theorem~\ref{fbl:thm:k1-triples} and the
partition optimization, the clique-fibre choice function and the Lindenbaum
recursion; no error. Independent programs confirmed the tables of
$A_2(n,1)$ and $A_3(n,1)$ for $n\le12$ from the closed forms, the critical
counts, $R_5=24$ by brute force, and $I_{C_4}$, $I_{C_5}$. The delivered
verifier, run on a copy, printed its recorded output. The write reread every
proof, found no error, reran the verifier on a copy (about 6~seconds,
output identical up to line endings), and made the computation reported in
Section~\ref{fbl:sec:further}.

\emph{Relation to the repository.}
\begin{itemize}\raggedright
\item \emph{Formal status.} No statement of this report is formalized in
 Lean or Rocq, and its place in the collection confers no formal status.
 The plan of Section~\ref{fbl:sec:formal} is a proposal; at the write's base
 commit no Lean or Rocq file of the repository mentions blacklists, free
 sets of set mappings, Freiling's axiom, $\SAP_k$ or the Boolean prime ideal theorem.
\item \emph{Stale claims.} The source's one statement about the repository,
 that it ``already contains substantial work on measurable box games and
 infinite hat guessing'', is correct at the pin and at the write's base
 commit: \file{measurable-box-games} (in this category, five Parts; title
 ``Countable Information, Uncountably Many Boxes'') is unchanged between the
 two. Nothing needed correcting. The bibliography entry~\cite{ProveItBoxes}
 names the repository's owner as its author; that report's README records
 how each of its Parts was produced.
\item \file{measurable-box-games} is cited as background only: it shares the
 motivation from Glazer's work and no theorem with this report.
\item \file{bounded-width-power-set-compression} (same category and batch;
 \emph{see also}): it too pairs an exact finite extremal theory with its
 choiceless infinite form over ZF (there, an infinite set admitting a
 bounded-width compression of its power set is Dedekind-infinite), and it
 too asks for the weakest standard choice principle that suffices, naming
 finite-choice schemes, the Boolean prime ideal theorem and compactness
 principles among the candidates, as Questions~\ref{fbl:q:sap1}
 and~\ref{fbl:q:hierarchy} do here. Neither source cites the other; no
 theorem is shared.
\item Batch~114 also filed \file{noisy-parity-cubes},
 \file{robust-neutral-choice} and \file{random-bits-arithmetic-cuts} in the
 same category; they share only the Glazer motivation and choice-built
 nonmeasurable objects, and no theorem.
\end{itemize}
\end{writenote}

\begin{writenote}[What is not claimed, collected from the source]
The partition of an infinite system into $2k+1$ safe sets is de Bruijn and
Erd\H{o}s's (1951), ``not \ldots\ this report'', and their paper assumes the
Axiom of Choice. ``No exhaustive priority search has been made for the
finite refinements proved here'', and the delivery README makes ``no
historical-priority claim \ldots\ without a more exhaustive literature
review''; the ledger of Section~\ref{fbl:sec:ledger} must be read the same
way (see the reading conventions below). No endorsement, collaboration or
authorship by Elliot Glazer is implied; the motivation rests on his public
work only. The computation ``is verification, not proof''. The perfect safe
set is ``classical tools, derived here'': Kuratowski--Ulam and Mycielski,
with the application to blacklists spelled out. The regularity theorems do
not yield a partition of $X$ into finitely many definable safe sets. The
coefficientwise conjectures are ``supported only by the proved coefficients
and small exhaustive checks'', and the exact choice strength of $\SAP_k$ is
open. The research program and the formalization plan are proposals.
\end{writenote}

\begin{writenote}[Reading conventions]
In the ledger of Section~\ref{fbl:sec:ledger}, ``no priority claim after an
exhaustive search'' means that no priority is claimed, as a claim would need
an exhaustive search; it does not say that one was made (the status box says
the opposite). ``Proved here'' in that ledger means proved in this text, not
first proved here. The colour set $2k+1$ in $c:X\to2k+1$ is the von Neumann
ordinal $\{0,\dots,2k\}$. ``Safe'' is de Bruijn and Erd\H{o}s's
``independent'' and the set-mapping literature's ``free''. The source reuses
letters as follows; no symbol was renamed.
\end{writenote}
\begin{center}\small
\renewcommand{\arraystretch}{1.12}
\begin{tabularx}{\textwidth}{@{}>{\raggedright\arraybackslash}p{0.17\textwidth}>{\raggedright\arraybackslash}X>{\raggedright\arraybackslash}p{0.27\textwidth}@{}}
\toprule
Symbol & Meaning here (where) & Not to be read as\\
\midrule
$F$, $F(x)$ & the blacklist system and the blacklist of $x$ & the families
 $\cA_e$, $\cA$\\
$A_r(n,k)$, $A$ & the minimum number of safe $r$-sets
 \eqref{fbl:eq:Ar-definition}; a safe set (Corollary~\ref{fbl:cor:weighted})
 or a member of the family $\cA$ (proof of
 Theorem~\ref{fbl:thm:choice-sandwich}) & each other\\
$\cA_e$, $\cA$ & the $r$-sets containing the edge $e$ (proof of
 Theorem~\ref{fbl:thm:asymptotic}); a family of $m$-element sets
 (Theorem~\ref{fbl:thm:choice-sandwich}) & each other\\
$\Cm_m$, $\Cm_{\mathrm{fin}}$ & choice for families of $m$-element sets,
 for families of nonempty finite sets (Section~\ref{fbl:sec:further}) &
 $C=R\cup R^{-1}$ (proof of Theorem~\ref{fbl:thm:perfect}); the cycle
 $C_s$; countable choice\\
$q$, $Q(H)$ & $q=2k+1$; the cost of a component $H$ (proof of
 Theorem~\ref{fbl:thm:k1-triples}) & each other\\
$t$, $t(G)$ & $t=r-1$, the number of blocks (Section~\ref{fbl:sec:critical});
 the number of triangles of $G$ & each other\\
$R$, $R_q$ & a relation on $X$ (Section~\ref{fbl:sec:definable}); the
 number of labelled regular tournaments on $q$ points & each other\\
$H_n$, $H$ & the conjectured minimizer
 (Conjecture~\ref{fbl:conj:pseudoforest-poly}); a unicyclic component &
 each other\\
$s$, $a$ & a block size (proof of Theorem~\ref{fbl:thm:alpha}); the
 remainder in $n=3a+s$ and the order of a component
 (Theorem~\ref{fbl:thm:k1-triples}); $a$ is also a point of $A$ in the
 sandwich proof & \\
$m$ & the number of edges in \eqref{fbl:eq:i3-formula}; the size of the
 sets in $\Cm_m$ & each other\\
$c$, $c_n$ & a colouring or a colour; the correction \eqref{fbl:eq:cn} & \\
$b(x)$, $L(x)$ & $\lvert F(x)\rvert$ and the list of $x$
 (Theorem~\ref{fbl:thm:list}) & \\
$e$, $e_F(S)$, $e(G)$ & a conflict edge; edge counts & \\
$i_r$, $I_G(z)$, $\alpha$ & the number of safe (independent) $r$-sets, the
 independence polynomial, the independence number & \\
$\SAP_k$, $\SAP_{\mathrm{fin}}$, $\SAP^{+}_{\mathrm{fin}}$, $\BPI$ & the
 partition principle for $k$-bounded systems; its analogues for finite
 blacklists with countably many colours (Section~\ref{fbl:sec:further}); the
 Boolean prime ideal theorem &
 \\
$P$, $\mu$ & a perfect set; a probability measure & \\
\bottomrule
\end{tabularx}
\end{center}

\section{The sparse symmetry-avoidance problem}

\begin{definition}[Blacklist system]
Let $X$ be a set and let $k\in\N$.  A \emph{$k$-bounded blacklist system} is
a map
\[
   F:X\longrightarrow [X]^{\le k}
\]
with $x\notin F(x)$ for every $x\in X$.  We draw an arc $x\to y$ when
$y\in F(x)$.
\end{definition}

\begin{definition}[Safe set and conflict graph]
Distinct $x,y\in X$ are \emph{compatible} when
\[
       y\notin F(x)\quad\text{and}\quad x\notin F(y).
\]
A subset $A\subseteq X$ is \emph{safe} when every two distinct members are
compatible.  The \emph{conflict graph} $G_F$ is the undirected graph on $X$
with
\[
   \{x,y\}\in E(G_F)
   \quad\Longleftrightarrow\quad
   y\in F(x)\ \text{or}\ x\in F(y).
\]
Thus safe sets are exactly the independent sets of $G_F$.
\end{definition}

For finite $X$, write $n=\lvert X\rvert$ and let
$i_r(F)=i_r(G_F)$ denote the number of safe $r$-subsets.  The central finite
spectrum is
\begin{equation}\label{fbl:eq:Ar-definition}
   A_r(n,k)
   \defeq
   \min_F i_r(F),
\end{equation}
where the minimum ranges over all $k$-bounded systems on an $n$-element
labelled set.

\begin{problem}[Sparse symmetry-avoidance spectrum]\label{fbl:prob:main}
Determine $A_r(n,k)$, the structure and number of its minimizers, and the
corresponding infinite partition principle.  In particular:
\begin{enumerate}[label=(\alph*)]
\item When is a safe $r$-set forced?
\item What happens at the last value of $n$ for which it can fail?
\item How many safe $r$-sets are forced above threshold?
\item Which weak choice principle is exactly needed to partition every
infinite system into $2k+1$ safe sets?
\item What stronger conclusions follow from Borel or measurable regularity?
\end{enumerate}
\end{problem}

The finite and infinite questions use the same object but behave very
differently.  Finite systems admit a deterministic elimination algorithm.
Infinite systems require a compactness step, and that step has genuine
foundational content.

\section{The graph translation and local sparsity}

For $S\subseteq X$, let $G_F[S]$ be the induced conflict graph and let
$e_F(S)$ be its number of edges.

\begin{lemma}[Induced edge budget]\label{fbl:lem:edge-budget}
For every finite $S\subseteq X$,
\begin{equation}\label{fbl:eq:edge-budget}
   e_F(S)
   \le
   \sum_{x\in S}\lvert F(x)\cap S\rvert
   \le k\lvert S\rvert.
\end{equation}
\end{lemma}

\begin{proof}
Every conflict edge in $G_F[S]$ is witnessed by at least one directed
blacklist arc whose tail and head both lie in $S$.  Counting such arcs gives
the first inequality; the blacklist bound gives the second.
\end{proof}

Consequently every nonempty finite induced subgraph has average degree at
most $2k$ and hence a vertex of degree at most $2k$.  A more useful local
version allows the budget to vary from point to point.

\begin{theorem}[Heterogeneous list theorem]\label{fbl:thm:list}
Let $X$ be finite, put $b(x)=\lvert F(x)\rvert$, and assign a finite list of
colours $L(x)$ to each $x\in X$.  If
\begin{equation}\label{fbl:eq:local-list}
          \lvert L(x)\rvert\ge 2b(x)+1
          \qquad(x\in X),
\end{equation}
then there is a proper list-colouring of $G_F$; equivalently, one may choose
$c(x)\in L(x)$ so that every colour class is safe.
\end{theorem}

\begin{proof}
For a nonempty $S\subseteq X$, suppose every $x\in S$ had
$d_{G_F[S]}(x)\ge 2b(x)+1$.  Then
\[
  2e_F(S)=\sum_{x\in S}d_{G_F[S]}(x)
  \ge 2\sum_{x\in S}b(x)+\lvert S\rvert,
\]
contradicting \cref{fbl:lem:edge-budget}.  Hence some $x\in S$ has
$d_{G_F[S]}(x)\le 2b(x)$.  Remove such a vertex repeatedly.  Colour the
vertices in reverse removal order.  When $x$ is restored, at most $2b(x)$
colours are forbidden by already coloured neighbours, while its list has at
least $2b(x)+1$ colours.
\end{proof}

\begin{corollary}[Finite partition bound]\label{fbl:cor:qcolor}
Every finite $k$-bounded blacklist system is the union of
\[
              q\defeq 2k+1
\]
safe sets.
\end{corollary}

\begin{corollary}[Weighted safe set]\label{fbl:cor:weighted}
If nonnegative weights $w(x)$ are assigned to a finite $k$-bounded system,
then some safe set $A$ satisfies
\[
       \sum_{x\in A}w(x)
       \ge \frac{1}{2k+1}\sum_{x\in X}w(x).
\]
\end{corollary}

\begin{proof}
Apply \cref{fbl:cor:qcolor} and take a colour class of maximum total weight.
\end{proof}

\begin{writenote}[Credit, 5 October 2026]
Corollary~\ref{fbl:cor:qcolor}, with the proof given here for constant
lists (at most $k\lvert S\rvert$ edges, so a vertex of degree at most $2k$,
then induction), is the finite case of Theorem~3 of de Bruijn and
Erd\H{o}s~\cite{deBruijnErdos} and is proved there the same way (note in
Section~\ref{fbl:sec:dbe-credit}). Their paper has no list
version: Theorem~\ref{fbl:thm:list}, with lists of size $2\lvert
F(x)\rvert+1$ varying from point to point, and
Corollary~\ref{fbl:cor:weighted} are not in it.
\end{writenote}

\section{The exact finite threshold}

The colouring bound immediately gives a safe set of size
$\lceil n/q\rceil$.  The next result shows that no improvement is possible,
for any value of $n$.

\begin{theorem}[Exact finite avoidance number]\label{fbl:thm:alpha}
For every $n\ge1$ and $k\ge0$,
\begin{equation}\label{fbl:eq:alpha}
  \min_F \alpha(G_F)
  =
  \left\lceil\frac{n}{2k+1}\right\rceil.
\end{equation}
\end{theorem}

\begin{proof}
The lower bound follows from \cref{fbl:cor:qcolor}: one of the $q=2k+1$ colour
classes has at least $\lceil n/q\rceil$ vertices.

For sharpness, partition the vertex set into blocks of size at most $q$, with
exactly $\lceil n/q\rceil$ blocks.  On a block of odd size $s=2h+1\le q$,
use the cyclic regular tournament, so every vertex has outdegree $h\le k$.
On a block of even size $s=2h\le q-1$, delete one vertex from a regular
tournament on $s+1$ vertices; every remaining outdegree is at most $h\le k$.
Let each point blacklist its out-neighbours in its block.  The conflict graph
is a disjoint union of cliques, one per block, so a safe set contains at most
one point from each block.  Its maximum size is therefore the number of
blocks.
\end{proof}

\begin{corollary}[Sharp threshold for a safe $r$-set]\label{fbl:cor:threshold}
For $r\ge2$, every $k$-bounded system on $n$ points has a safe $r$-set if
and only if
\begin{equation}\label{fbl:eq:threshold}
       n>(r-1)(2k+1).
\end{equation}
Equivalently, the least forcing value is
\[
       n_{\min}(r,k)=(r-1)(2k+1)+1.
\]
\end{corollary}

This exact threshold is the first answer to
\cref{fbl:prob:main}.  Its equality case is much more rigid than the proof of the
lower bound suggests.

\section{Critical rigidity and exact enumeration}\label{fbl:sec:critical}

Put $q=2k+1$ and $t=r-1$.  The critical size is $n=tq$.

\begin{theorem}[Classification at the last failing value]\label{fbl:thm:critical}
Let $F$ be a $k$-bounded blacklist system on $n=tq$ points, where
$q=2k+1$ and $t\ge1$.  The following are equivalent.
\begin{enumerate}[label=(\roman*)]
\item $F$ has no safe $(t+1)$-set.
\item The conflict graph is a disjoint union of $t$ copies of $K_q$.
\item The vertices can be partitioned into $t$ blocks of size $q$ such that,
on each block, the blacklist arcs form a regular tournament and there are no
arcs between distinct blocks.
\end{enumerate}
\end{theorem}

\begin{proof}
The implications (iii)$\Rightarrow$(ii)$\Rightarrow$(i) are immediate.
Assume (i) and write $G=G_F$.  The complement $\overline G$ contains no
$K_{t+1}$, because such a clique would be a safe $(t+1)$-set.  Tur\'an's
theorem, including its equality case, gives
\[
 e(\overline G)
 \le \binom{t}{2}q^2.
\]
Therefore
\begin{align*}
 e(G)
 &\ge \binom{tq}{2}-\binom{t}{2}q^2\\
 &=\frac{tq(q-1)}{2}
 =ktq
 =kn.
\end{align*}
On the other hand, \cref{fbl:lem:edge-budget} gives $e(G)\le kn$.  Equality
holds throughout.  Equality in Tur\'an's theorem forces $\overline G$ to be
the balanced complete $t$-partite graph, hence $G$ is $tK_q$.

There are at most $kn$ blacklist arcs, while the $kn$ conflict edges each
need at least one witnessing arc.  Thus there is exactly one arc on every
conflict edge, every point uses all $k$ blacklist slots, and no pair is
blacklisted in both directions.  Each $K_q$ is therefore a tournament with
constant outdegree $k$, that is, a regular tournament.
\end{proof}

The theorem explains why the odd number $2k+1$ is not merely a consequence
of averaging: at equality it is the unique block size supporting a complete
conflict graph with exactly $k$ outgoing arcs at every point.

Let $R_q$ be the number of regular tournaments on a fixed labelled
$q$-element vertex set.  Regular tournaments exist because $q$ is odd.

\begin{corollary}[Number of labelled critical systems]\label{fbl:cor:critical-count}
The number of labelled $k$-bounded systems on $n=tq$ points with no safe
$(t+1)$-set is
\begin{equation}\label{fbl:eq:critical-count}
  \boxed{
  \frac{(tq)!}{(q!)^t t!}\,R_q^t }.
\end{equation}
\end{corollary}

\begin{proof}
Choose an unordered partition into $t$ blocks of size $q$, then choose one of
the $R_q$ regular tournaments independently on each block.
\end{proof}

For example, $R_3=2$ and $R_5=24$.  Thus there are
\[
 \frac{6!}{(3!)^2 2!}\,2^2=40
\]
critical systems for $(k,r,n)=(1,3,6)$, and
\[
 \frac{10!}{(5!)^2 2!}\,24^2=72{,}576
\]
critical systems for $(k,r,n)=(2,3,10)$.

\section{The safe-subset spectrum}

The threshold asks only whether $A_r(n,k)$ vanishes.  The more informative
question is how rapidly it grows after the threshold.

\subsection{Safe pairs: an exact formula}

\begin{theorem}[Exact pair spectrum]\label{fbl:thm:pairs}
For all $n\ge1$ and $k\ge0$,
\begin{equation}\label{fbl:eq:pairs}
   \boxed{
   A_2(n,k)=\max\left\{0,\binom n2-kn\right\}. }
\end{equation}
\end{theorem}

\begin{proof}
There are at most $kn$ conflict edges, so at least
$\binom n2-kn$ pairs are safe.  The count is nonnegative.

If $n\le2k+1$, orient the complete graph with maximum outdegree at most $k$
as in the proof of \cref{fbl:thm:alpha}; there are no safe pairs.  If
$n>2k+1$, label the vertices by $\Z/n\Z$ and set
\[
       F(i)=\{i+1,i+2,\ldots,i+k\}\pmod n.
\]
No unordered pair receives two opposite arcs, so the conflict graph has
exactly $kn$ edges and the lower bound is attained.
\end{proof}

\subsection{All fixed orders: a sharp two-term asymptotic}
We use the convention $\binom ab=0$ when $b<0$ or $b>a$.

\begin{theorem}[Universal lower bound and asymptotic sharpness]
\label{fbl:thm:asymptotic}
For every $n,k,r$,
\begin{equation}\label{fbl:eq:union-lower}
 A_r(n,k)
 \ge
 \max\left\{0,
 \binom nr-kn\binom{n-2}{r-2}
 \right\}.
\end{equation}
If $n>2k+1$, then
\begin{align}
 A_r(n,k)
 \le{}&
 \binom nr-kn\binom{n-2}{r-2}
 \nonumber\\
 &+n\binom{2k}{2}\binom{n-3}{r-3}
 +\binom{kn}{2}\binom{n-4}{r-4}.
 \label{fbl:eq:bonf-upper}
\end{align}
Consequently, for fixed $k$ and $r\ge2$,
\begin{equation}\label{fbl:eq:asymptotic}
 A_r(n,k)
 =
 \binom nr-kn\binom{n-2}{r-2}
 +O_{k,r}(n^{r-2}),
\end{equation}
and hence
\begin{equation}\label{fbl:eq:density-asymptotic}
 \frac{A_r(n,k)}{\binom nr}
 =1-\frac{k r(r-1)}{n-1}+O_{k,r}(n^{-2}).
\end{equation}
\end{theorem}

\begin{proof}
For each conflict edge $e$, let $\cA_e$ be the family of $r$-subsets
containing $e$.  Then $\lvert\cA_e\rvert=\binom{n-2}{r-2}$.  Since there are
at most $kn$ conflict edges, the union bound shows that at most
$kn\binom{n-2}{r-2}$ subsets are unsafe, proving
\cref{fbl:eq:union-lower}.

For the upper bound on the minimum, use the circulant construction from
\cref{fbl:thm:pairs}.  Its conflict graph is $2k$-regular with $kn$ edges.
Bonferroni's inequality gives
\[
 \left|\bigcup_e\cA_e\right|
 \ge
 \sum_e|\cA_e|
 -\sum_{e<f}|\cA_e\cap\cA_f|.
\]
There are exactly $n\binom{2k}{2}$ pairs of incident edges, each lying in
$\binom{n-3}{r-3}$ common $r$-sets.  The number of disjoint edge pairs is at
most $\binom{kn}{2}$, and each lies in $\binom{n-4}{r-4}$ common $r$-sets.
Subtracting the resulting lower bound on unsafe sets from $\binom nr$ proves
\cref{fbl:eq:bonf-upper}.  The two error terms are
$O_{k,r}(n^{r-2})$, and division by $\binom nr$ gives
\cref{fbl:eq:density-asymptotic}.
\end{proof}

The first correction has a transparent interpretation: to first order, each
of the $kn$ possible conflict edges destroys
$\binom{n-2}{r-2}$ candidate safe sets.  Overlaps among destroyed families
enter one order later.

\subsection{An exact third coefficient when \texorpdfstring{$k=1$}{k=1}}
For $k=1$, the conflict graphs are exactly the finite pseudoforests: every
connected component contains at most one cycle.  This permits an exact safe
triple count.

\begin{theorem}[Exact safe triples for one blacklist per point]
\label{fbl:thm:k1-triples}
For $n<3$, $A_3(n,1)=0$.  For $n\ge3$, write $n=3a+s$ with
$s\in\{0,1,2\}$.  Then
\begin{equation}\label{fbl:eq:k1-triples}
 \boxed{
 A_3(n,1)
 =\binom n3-n(n-2)+c_n, }
\end{equation}
where
\begin{equation}\label{fbl:eq:cn}
 c_n=
 \begin{cases}
    2a,   &s=0,\\
    2a+2, &s=1,\\
    2a+3, &s=2.
 \end{cases}
\end{equation}
The bound is attained by a disjoint union of triangles when $s=0$, by
$(a-1)K_3\sqcup C_4$ when $s=1$, and by
$(a-1)K_3\sqcup C_5$ when $s=2$.
\end{theorem}

\begin{proof}
A graph is the conflict graph of a $1$-bounded system exactly when it is a
pseudoforest.  One direction follows from
$e(S)\le|S|$.  Conversely, orient the unique cycle of each unicyclic
component cyclically and every attached tree toward the cycle; orient a tree
toward an arbitrary root.  Every outdegree is at most one.

For any graph $G$ with $m$ edges, degrees $d(v)$, and $t(G)$ triangles,
inclusion--exclusion on the edges contained in a triple gives
\begin{equation}\label{fbl:eq:i3-formula}
 i_3(G)
 =\binom n3-m(n-2)
   +\sum_v\binom{d(v)}2-t(G).
\end{equation}
Adding an edge without leaving the class of pseudoforests cannot increase
the number of independent triples.  Hence a minimizer may be assumed to have
all components unicyclic, and therefore $m=n$.

For a unicyclic component $H$ of order $s$, set
\[
 Q(H)=\sum_{v\in H}\binom{d(v)}2-t(H).
\]
If $H=C_3$, then $Q(H)=2$.  If $s\ge4$ and all degrees equal two, then
$H=C_s$ and $Q(H)=s$.  Otherwise the degree sum is $2s$ but the degrees are
not all two, so integer convexity gives
$\sum_v\binom{d(v)}2\ge s+1$.  A unicyclic component has at most one
triangle, whence again $Q(H)\ge s$.

Thus a component of size three costs at least $2$, while every other
component of size $s\ge4$ costs at least $s$.  Minimizing the total cost over
partitions of $n$ into integers at least three gives exactly
\cref{fbl:eq:cn}: use as many $3$'s as possible, with one $4$ or one $5$ to
absorb a nonzero remainder.  Substitution into \cref{fbl:eq:i3-formula} proves
the result, and the displayed cycle unions attain equality.
\end{proof}

\begin{writenote}[A step made explicit, 5 October 2026]
The step ``a minimizer may be assumed to have all components unicyclic''
needs that a pseudoforest on $n\ge3$ vertices to which no edge can be added
without leaving the class has only unicyclic components. Suppose such a
pseudoforest has a tree component $T$. If $\lvert T\rvert\ge3$, two
nonadjacent vertices of $T$ exist, and joining them makes $T$ unicyclic. If
$\lvert T\rvert\le2$, then $n\ge3$ gives another component $U$, and an edge
from $T$ to $U$ yields a tree if $U$ is a tree and a unicyclic component if
$U$ is unicyclic. Either way an edge can be added, a contradiction. Adding
edges never increases $i_3$, so starting from any minimizer and adding
edges while possible reaches a minimizer with all components unicyclic,
hence with $m=n$. The minimum over partitions is as stated: a part of size
$3$ costs $2=3-1$ and a part of size $s\ge4$ costs at least $s$, so the
total cost is at least $n$ minus the number of parts of size~$3$, which is
at most $a$ when $s=0$ and at most $a-1$ otherwise (the remainder must be
absorbed by a part of size at least $4$); this gives $2a$, $2a+2$, $2a+3$.
\end{writenote}

The first values are:
\begin{center}
\begin{tabular}{c|rrrrrrrrrrrr}
$n$&1&2&3&4&5&6&7&8&9&10&11&12\\ \hline
$A_2(n,1)$&0&0&0&2&5&9&14&20&27&35&44&54\\
$A_3(n,1)$&0&0&0&0&0&0&6&15&27&48&75&108
\end{tabular}
\end{center}

\section{Algorithms and certificates}\label{fbl:sec:algorithms}

The proof of \cref{fbl:thm:list} is constructive.

\begin{proposition}[Linear-time peeling after graph construction]
\label{fbl:prop:algorithm}
Given the blacklist lists explicitly, one can produce the elimination order
and a safe $(2k+1)$-colouring in
$O(n+\sum_x|F(x)|+e(G_F))$ time, using adjacency lists and a queue of
vertices satisfying $d(x)\le2|F(x)|$.
\end{proposition}

\begin{proof}[Algorithm]
Construct the undirected conflict graph while deduplicating opposite arcs.
Maintain current degrees and enqueue every currently eligible vertex.
Removing a vertex can only decrease the degrees of its neighbours, so each
vertex and edge is processed a constant number of times.  Reverse the
removal order and greedily choose an available colour (or an available list
colour in the heterogeneous version).
\end{proof}

This yields three useful certificates.
\begin{itemize}
\item A colouring certifies a safe set of size at least
$\lceil n/(2k+1)\rceil$ by its largest colour class.
\item At the critical size $n=(r-1)(2k+1)$, failure of a safe $r$-set can be
verified simply by checking the block decomposition and regular tournament
conditions of \cref{fbl:thm:critical}.
\item A claimed extremizer can be checked in time linear in its explicit
blacklist data, rather than by searching all $r$-subsets.
\end{itemize}

The optimization problem above threshold remains nontrivial: finding a
maximum independent set is hard in broad sparse graph classes, even though
the universal guarantee and its certificate are easy.

\begin{writenote}[Scope of the third certificate, 5 October 2026]
The text proves a linear-time check for two kinds of extremizer only: the
critical systems, through the block and regular-tournament conditions of
Theorem~\ref{fbl:thm:critical} (second item), and the minimizers of
$A_2(n,k)$, since $i_2(F)=\binom n2-e(G_F)$ and the proof of
Theorem~\ref{fbl:thm:pairs} shows that a system is a minimizer exactly when
$e(G_F)=\min\{kn,\binom n2\}$. For a claimed minimizer of $A_r(n,k)$ with
$r\ge3$, or of $\alpha(G_F)$ above the critical size, the text gives no such
check: verifying it needs the value $A_r(n,k)$, unknown in general
(Question~\ref{fbl:q:exact}), or an upper bound on the independence number.
The third item is therefore stated in Section~\ref{fbl:sec:further} as an
open question for those cases (item~\ref{fbl:fq:certificate}). The
hardness sentence above is the source's, given without a reference.
\end{writenote}

\section{The infinite partition and its choice content}

\subsection{The classical compactness extension}

\begin{theorem}[de Bruijn--Erd\H{o}s, relation form]\label{fbl:thm:dbe}
Assume a graph-colouring compactness principle, for example the Boolean
Prime Ideal Theorem.  If $F:X\to[X]^{\le k}$, then $X$ is the union of
$2k+1$ safe sets.
\end{theorem}

\begin{proof}
Every finite induced subsystem is $(2k+1)$-colourable by
\cref{fbl:cor:qcolor}.  Graph-colouring compactness extends the compatible
finite constraints to a colouring of all of $X$.  This is the relation
problem proved by de Bruijn and Erd\H{o}s \cite{deBruijnErdos}; the weaker
hypothesis $\BPI$ suffices because it is equivalent to propositional
compactness \cite{Cowen1990,Cowen1998,HowardRubin}.
\end{proof}

The finite proof does not directly recurse through an infinite set.  An
infinite conflict graph may have no vertex of degree at most $2k$, because a
point can be blacklisted by arbitrarily many others.  Compactness is doing
real work.

\subsection{A precise ZF principle}

\begin{definition}[Sparse Avoidance Partition principle]
For $k\ge1$, $\SAP_k$ is the statement:
\begin{quote}
For every set $X$ and every $F:X\to[X]^{\le k}$ with $x\notin F(x)$, there
exists $c:X\to 2k+1$ such that every colour class is safe.
\end{quote}
\end{definition}

\begin{theorem}[Choice sandwich]\label{fbl:thm:choice-sandwich}
Over ZF, for every $k\ge1$,
\begin{equation}\label{fbl:eq:choice-sandwich}
 \BPI
 \Longrightarrow \SAP_k
 \Longrightarrow
 \bigwedge_{m=2}^{k+1}\Cm_m,
\end{equation}
where $\Cm_m$ denotes the axiom of choice for families of $m$-element sets.
\end{theorem}

\begin{proof}
The first implication is \cref{fbl:thm:dbe}.  For the second, fix
$2\le m\le k+1$ and a family $\cA$ of pairwise arbitrary $m$-element sets.
Disjointify it by taking
\[
 X=\{(A,a):A\in\cA,\ a\in A\}.
\]
For $x=(A,a)$, let
\[
 F(x)=\{(A,b):b\in A\setminus\{a\}\}.
\]
The blacklist size is $m-1\le k$, and each fibre
$\{A\}\times A$ induces a clique.  A $\SAP_k$ colouring therefore assigns
distinct colours to the $m$ points in every fibre.  Choose from $A$ the
unique element receiving the least colour used on that fibre.  This defines
a choice function for $\cA$.
\end{proof}

Since choice for pairs is not provable in ZF, neither is $\SAP_k$ for any
$k\ge1$.  On the other hand, \cref{fbl:eq:choice-sandwich} leaves a large gap:
$\BPI$ is substantially stronger than the displayed bounded finite-choice
principles.  Locating $\SAP_k$ in that gap is one of the main open problems
proposed here.

\begin{writenote}[5 October 2026]
For $k=1$ the lower end can be raised: $\SAP_1$ also implies choice for
triples (Proposition~\ref{fbl:prop:three}). The finite-blacklist analogue,
de Bruijn and Erd\H{o}s's countable-colour theorem, lies over ZF between
$\BPI$ and choice for families of nonempty finite sets
(Proposition~\ref{fbl:prop:finite}). The comparison of $\BPI$ with the
finite-choice principles in the preceding paragraph is stated without proof
or reference (item~\ref{fbl:fq:hierarchy} of
Section~\ref{fbl:sec:further}).
\end{writenote}

\subsection{Well-orderable sets require no additional choice}

\begin{theorem}[ZF theorem for well-orderable domains]\label{fbl:thm:wellorderable}
In ZF, if $X$ is well-orderable and $F:X\to[X]^{\le k}$, then $X$ has a
safe $(2k+1)$-colouring.
\end{theorem}

\begin{proof}
Introduce propositional atoms $p_{x,c}$ for $x\in X$ and
$c<2k+1$, and impose the finite clauses saying that each point receives
exactly one colour and conflict neighbours receive different colours.
Every finite set of clauses is satisfiable by \cref{fbl:cor:qcolor}.

Because $X$ is well-orderable, the set of all formulas in this language is
well-orderable.  The usual Lindenbaum construction can therefore be carried
out by transfinite recursion in ZF: at each formula $\varphi$, add either
$\varphi$ or $\neg\varphi$, choosing the first option when it preserves
finite satisfiability and the second otherwise.  At least one option
preserves finite satisfiability, since failure of both would be witnessed by
two finite contradictory subtheories.  The resulting complete finitely
satisfiable theory determines a truth assignment satisfying all original
clauses, hence the desired colouring.
\end{proof}

Thus any choiceless counterexample to $\SAP_k$ must live on a genuinely
non-well-orderable set.  This is exactly the kind of boundary between local
finite mathematics and global choice that motivates Glazer's work.

\section{Definable infinite relations: perfect and random safe sets}\label{fbl:sec:definable}

The unrestricted partition problem is choice-sensitive.  Borel or measurable
regularity produces a different and much stronger kind of conclusion: one
can find a very large single safe set even when blacklists are merely
countable.

\subsection{A perfect safe set from Baire category}

\begin{theorem}[Perfect symmetry avoidance]\label{fbl:thm:perfect}
Let $X$ be a perfect Polish space and let $R\subseteq X^2$ be a Borel
irreflexive relation whose vertical sections
\[
       R_x=\{y:(x,y)\in R\}
\]
are countable.  Then there is a nonempty perfect set $P\subseteq X$ such
that for all distinct $x,y\in P$,
\[
        (x,y)\notin R
        \quad\text{and}\quad
        (y,x)\notin R.
\]
In blacklist language, $P$ is a perfect safe set.
\end{theorem}

\begin{proof}
Every countable subset of a perfect Polish space is meagre.  By the
Kuratowski--Ulam theorem, the Borel set $R$ is meagre in $X^2$, because all
of its vertical sections are meagre.  Coordinate exchange is a homeomorphism,
so $R^{-1}$ is meagre as well.  The symmetric conflict relation
$C=R\cup R^{-1}$ is therefore meagre.  Mycielski's theorem supplies a
perfect set $P$ whose distinct pairs avoid $C$ \cite{Mycielski,Kechris}.
\end{proof}

\begin{corollary}[Definable strengthening of Freiling symmetry]\label{fbl:cor:freiling}
A Borel countable-valued blacklist relation on $\R$ not only has a safe
pair; it has a perfect, hence continuum-sized, safe set.
\end{corollary}

This theorem cleanly separates regular assignments from the classical
counterexample under CH.  If the reals are well-ordered in type $\omega_1$,
letting each real blacklist all of its predecessors gives countable
blacklists and destroys every safe pair.  Such a relation cannot be Borel.

\subsection{An almost surely safe countable sample}

\begin{theorem}[Random countable avoidance]\label{fbl:thm:random}
Let $(X,\mu)$ be an atomless standard probability space and let
$R\subseteq X^2$ be product-measurable with every vertical section $R_x$
countable.  If $(X_n)_{n\in\N}$ are independent samples with law $\mu$, then
with probability one:
\begin{enumerate}[label=(\roman*)]
\item the points $X_n$ are all distinct; and
\item for every $i\ne j$, neither $(X_i,X_j)$ nor $(X_j,X_i)$ belongs to
$R$.
\end{enumerate}
Thus the random range $\{X_n:n\in\N\}$ is almost surely a countably infinite
safe set.
\end{theorem}

\begin{proof}
Atomlessness makes every countable set null.  Tonelli's theorem gives
\[
  (\mu\times\mu)(R)
  =\int_X \mu(R_x)\,d\mu(x)=0.
\]
The inverse relation is also null, by coordinate symmetry, and the diagonal
is null.  For each ordered pair $i\ne j$, the event that
$(X_i,X_j)$ lies in $R\cup R^{-1}$ has probability zero; for each
$i<j$, the event $X_i=X_j$ has probability zero.  There are only countably
many such events, so their union is null.
\end{proof}

For every fixed $r$, an equivalent finite-dimensional statement is that
$\mu^r$-almost every $r$-tuple of distinct points is safe.  The countable
version makes the regularity amplification more vivid: measurable
countable blacklists are almost surely avoided simultaneously by an entire
random sequence.

\subsection{What these regularity theorems do not say}
Neither \cref{fbl:thm:perfect,fbl:thm:random} automatically yields a partition of
all of $X$ into finitely many definable safe sets.  Perfect-set extraction,
almost-sure sampling, ordinary finite colouring, Borel colouring, and
measurable colouring are different problems.  Their separation is a major
source of further questions in \cref{fbl:sec:questions}.

\section{Finite versus infinite: the phase diagram}

The same blacklist relation has four sharply different regimes.

\begin{center}
\begin{tabularx}{0.96\textwidth}{>{\bfseries}p{0.20\textwidth}p{0.22\textwidth}X}
Regime & Hypothesis & Best conclusion established here\\ \toprule
Finite, arbitrary & $|X|=n$, $|F(x)|\le k$ &
Exact safe number $\lceil n/(2k+1)\rceil$; exact threshold; critical rigidity;
quantitative spectrum.\\
Infinite, arbitrary & no regularity &
A $(2k+1)$-partition under $\BPI$; in ZF for well-orderable $X$; exact weak
choice strength open.\\
Topological & perfect Polish $X$, Borel countable sections &
A perfect safe set.\\
Measure-theoretic & atomless standard probability, measurable countable
sections & Almost every i.i.d. countable sample is safe.\\
\bottomrule
\end{tabularx}
\end{center}

This is the central conceptual answer to the design request.  The finite
form is exact and algorithmic.  The unrestricted infinite form is governed
by compactness and weak choice.  The definable infinite form is dramatically
more regular than either: category and measure amplify a pairwise symmetry
intuition into perfect or countably infinite safe sets.

\section{Computational verification}\label{fbl:sec:computation}

The accompanying script \texttt{verification.py} performs exhaustive checks
for all $1$-bounded systems on at most seven labelled vertices.  A vertex has
$n$ choices (no blacklist target, or one of the other $n-1$ vertices), so the
raw search contains $n^n$ systems.  Systems are deduplicated by conflict
graph before independent-set counting.

The executed output is:
\begin{verbatim}
Exhaustive minima for k=1
n  distinct conflict graphs  A_2  A_3  alpha_min
1                         1    0    0          1
2                         2    0    0          1
3                         8    0    0          1
4                        57    2    0          2
5                       608    5    0          2
6                      8524    9    0          2
7                    145800   14    6          3

Critical systems at (k,r,n)=(1,3,6): 40
Labeled regular tournaments: R_3=2, R_5=24
Critical systems at (k,r,n)=(2,3,10): 72576

All checks passed.
\end{verbatim}

The computation is verification, not proof.  The proofs in
\cref{fbl:thm:alpha,fbl:thm:critical,fbl:thm:pairs,fbl:thm:asymptotic,fbl:thm:k1-triples} are
independent of the search.

\begin{writenote}[What the program computes, 5 October 2026]
The shipped \file{code/verification.py} (standard library only) does four
things. (1)~For each $n\le7$ it runs through all $n^n$ one-bounded systems,
collects their distinct conflict graphs (the second column, which is
A133686 of the OEIS, labelled graphs whose components have at most one
cycle, as the proof of Theorem~\ref{fbl:thm:k1-triples} predicts), and
computes the coefficientwise minimum of their independence polynomials; it
asserts that the coefficients of $z^2$ and $z^3$ equal
\eqref{fbl:eq:pairs} and \eqref{fbl:eq:k1-triples} with $k=1$, and that
the largest $r$ with a positive minimum is $\lceil n/3\rceil$. The column
headed \texttt{alpha\_min} prints that theorem value after the assertion.
(2)~It counts the systems on six points with no safe triple among all
$6^6=46{,}656$ one-bounded systems: $40$, by enumeration. (3)~It finds
$R_3=2$ and $R_5=24$ by running through all orientations of $K_3$ and
$K_5$. (4)~It obtains $72{,}576$ by evaluating
\eqref{fbl:eq:critical-count} with $R_5=24$; no two-bounded system is
enumerated, so this count rests on Corollary~\ref{fbl:cor:critical-count}.
The program writes no file. The write reran it on a copy (Windows, Python
3.14.4): about 6~seconds, output identical to
\file{data/verification_output.txt} up to line endings. For $n\le7$ every
coefficient of $I_{H_n}$ beyond $z^3$ vanishes, so these checks do cover
every coefficient of Conjecture~\ref{fbl:conj:pseudoforest-poly} there, as
the text after it says; Proposition~\ref{fbl:prop:small} turns this into a
proof for $n\le9$.
\end{writenote}

\section{A proposed research program}\label{fbl:sec:questions}

\subsection{The choice-theoretic core}

\begin{question}[Exact strength of $\SAP_1$]\label{fbl:q:sap1}
Is $\SAP_1$ equivalent over ZF to choice for pairs, to a known permutation
or multiple-choice principle, or to a strictly stronger principle?
\end{question}

The directed structure for $k=1$ is a partial functional graph.  Its
components may have cycles, one-way tails, or highly branching inverse
trees.  Colouring all components simultaneously appears to require choosing
phases or markers, but much less than arbitrary propositional compactness.

\begin{question}[The hierarchy $\SAP_k$]\label{fbl:q:hierarchy}
Do the principles $\SAP_k$ form a strict hierarchy?  Does some finite $k$
reach $\BPI$?  What are their implications among the standard principles in
Howard--Rubin's catalogue \cite{HowardRubin}?
\end{question}

\begin{question}[Choice-free countable-colour theorem]\label{fbl:q:countable}
De Bruijn and Erd\H{o}s also proved that finite, but not uniformly bounded,
blacklists admit a countable safe partition under their ambient choice
assumptions.  What is the exact ZF strength of that statement?
\end{question}

\subsection{The finite spectrum}

\begin{question}[Exact $A_r(n,k)$]\label{fbl:q:exact}
Determine $A_r(n,k)$ beyond pairs and the case $(r,k)=(3,1)$.  The asymptotic
in \cref{fbl:thm:asymptotic} fixes the first two orders but leaves an
$O(n^{r-2})$ optimization problem involving overlap patterns among conflict
edges.
\end{question}

For a graph $G$, let
\[
       I_G(z)=\sum_{r\ge0}i_r(G)z^r
\]
be its independence polynomial.

\begin{conjecture}[Coefficientwise pseudoforest minimizer]
\label{fbl:conj:pseudoforest-poly}
For $n\ge3$, among all $n$-vertex pseudoforests, the independence polynomial
is coefficientwise minimized by
\[
 H_n=
 \begin{cases}
  aK_3, &n=3a,\\
  (a-1)K_3\sqcup C_4, &n=3a+1,\\
  (a-1)K_3\sqcup C_5, &n=3a+2.
 \end{cases}
\]
Equivalently,
\[
 I_G(z)\succeq
 \begin{cases}
  (1+3z)^a,\\
  (1+3z)^{a-1}(1+4z+2z^2),\\
  (1+3z)^{a-1}(1+5z+5z^2),
 \end{cases}
\]
according to $n\bmod 3$.
\end{conjecture}

The theorem on safe triples proves one coefficient of this conjecture, and
the exhaustive search verifies all coefficients for $n\le7$.

\begin{writenote}[5 October 2026]
Theorem~\ref{fbl:thm:pairs} with $k=1$ proves the coefficient of $z^2$ as
well. Proposition~\ref{fbl:prop:small} proves both conjectures in their
smallest cases (this one for $n\le9$), and an exact computation of the write
confirms this one for $n\le18$ (item~\ref{fbl:fq:pseudoforest} of
Section~\ref{fbl:sec:further}).
\end{writenote}

\begin{conjecture}[Critical block minimization]\label{fbl:conj:critical-block}
When $q=2k+1$ divides $n$, the disjoint union $(n/q)K_q$ minimizes the
independence polynomial coefficientwise among conflict graphs of
$k$-bounded systems.
\end{conjecture}

The divisibility hypothesis matters.  For example, with $k=1$ and $n=4$,
$C_4$ has fewer independent pairs than $K_3\sqcup K_1$.

\begin{question}[Stability and supersaturation]\label{fbl:q:stability}
If a critical or near-critical system has very few safe $r$-sets, must it be
close, in edit distance, to a union of regular-tournament blocks?  Find sharp
supersaturation inequalities above $n=(r-1)(2k+1)$.
\end{question}

\begin{question}[Enumeration]\label{fbl:q:enumeration}
Enumerate all minimizers of $A_r(n,k)$ and obtain asymptotics for their
number.  At the exact threshold the answer is
\cref{fbl:eq:critical-count}; beyond it, one expects a mixture of tournament
counts and sparse stability theory.
\end{question}

\subsection{Definability and regularity}

\begin{question}[Borel partition number]\label{fbl:q:borel}
For a Borel $k$-bounded blacklist relation on a standard Borel space, is a
Borel safe partition into finitely many pieces always possible?  If not,
what is the optimal Borel chromatic bound, and how does it depend on
additional bounds on incoming sections?
\end{question}

\begin{question}[Measurable positive safe sets]\label{fbl:q:measurable}
Under the hypotheses of \cref{fbl:thm:random}, when must there exist a measurable
safe set of positive measure?  What is the best universal measure bound
under measure-preserving, bounded-indegree, or graphing assumptions?
\end{question}

\begin{question}[Regularity strength in ZF]\label{fbl:q:regularity}
Which fragments of dependent choice or Baire-category regularity suffice for
\cref{fbl:thm:perfect,fbl:thm:random}?  Can their exact foundational strengths be
separated from $\SAP_k$?
\end{question}

\subsection{Cardinal, probabilistic, and higher-arity extensions}

\begin{question}[Cardinal blacklists]\label{fbl:q:cardinal}
For infinite cardinals $\kappa$, study assignments
$F:X\to[X]^{\le\kappa}$.  Determine the least number of safe classes, the
largest forced safe cardinality, and the relation to generalized Freiling
principles and cardinal arithmetic.
\end{question}

\begin{question}[Random blacklist systems]\label{fbl:q:random}
In a random $k$-out directed model, determine the typical independence
number, safe-set polynomial, fluctuations, and the gap between typical and
worst-case behaviour.  The worst-case value is $n/(2k+1)$, while random
sparse graphs should have substantially larger safe sets.
\end{question}

\begin{question}[Hypergraph avoidance]\label{fbl:q:hypergraph}
Allow each point to forbid configurations of arity $d>2$.  Find the finite
threshold, compactness strength, and measurable perfect-set analogues for
safe subsets avoiding all directed forbidden patterns.
\end{question}

\begin{question}[Game and online forms]\label{fbl:q:game}
Study a game in which an adversary reveals blacklists online and a player
must colour immediately, or chooses a large safe set under limited queries.
Determine competitive ratios and whether randomization changes the sharp
constant $2k+1$.
\end{question}

\subsection{Sequence and formalization projects}
The arrays $A_r(n,k)$, the number of minimizers, and the labelled critical
counts \cref{fbl:eq:critical-count} are natural candidates for an OEIS-oriented
subproject once enough terms and independent literature checks are available.
The regular tournament factor $R_{2k+1}$ links the problem to existing
tournament enumeration, while the new two-parameter structure records the
finite--infinite symmetry problem rather than tournament counts alone.

\begin{writenote}[5 October 2026]
$R_{2k+1}$ is A007079 of the OEIS~\cite{OEIS} (labelled regular tournaments
on $2k+1$ nodes; $R_7=2640$). Nothing has been submitted to the OEIS. The
computation reported in Section~\ref{fbl:sec:further}
(item~\ref{fbl:fq:pseudoforest}) gives $A_r(n,1)$ for every $r$ and every
$n\le18$.
\end{writenote}

\subsection{Further questions and research}\label{fbl:sec:further}
\begin{writenote}[Added 5 October 2026, batch~114]
Following Vladimir Reshetnikov's standing rule of 4~October 2026, every
claim of the source that is not proved in full is stated here as an open
question, with its source, the argument offered and what is missing. The two
conjectures and the thirteen research questions above are the source's own
proposals; all are open, three are re-scoped below by
Propositions~\ref{fbl:prop:small}--\ref{fbl:prop:finite}, which the write
proves, and one by an exact computation. Four sentences of the body that
assert more than is proved are attached to the items they belong to: the
third certificate of Section~\ref{fbl:sec:algorithms} (a new item), the
comparison of $\BPI$ with the finite-choice principles after
Theorem~\ref{fbl:thm:choice-sandwich}, the sketch after
Question~\ref{fbl:q:sap1}, and the expectations in
Questions~\ref{fbl:q:enumeration} and~\ref{fbl:q:random}. The intake and the
write found no false statement in the source.
\end{writenote}
\begin{enumerate}[leftmargin=*,label=\textbf{\arabic*.},ref=\arabic*]
\item\label{fbl:fq:pseudoforest} \textbf{Coefficientwise pseudoforest
 minimizer} (Conjecture~\ref{fbl:conj:pseudoforest-poly}). Proved: the
 coefficients of $z^0,\dots,z^3$ (Theorems~\ref{fbl:thm:pairs}
 and~\ref{fbl:thm:k1-triples}), and with them the whole conjecture for
 $n\le9$ (Proposition~\ref{fbl:prop:small}). \emph{Observation of the
 write} (exact computation, 5~October 2026; program not shipped): the
 conjecture holds for every $n\le18$, so that $A_r(n,1)=[z^r]I_{H_n}$ for
 all $r$ and $n\le18$. Method: by the reduction in the proof of
 Theorem~\ref{fbl:thm:k1-triples} (every pseudoforest on $n\ge3$ vertices
 is a spanning subgraph of one whose components are all unicyclic, and
 adding edges lowers every coefficient), it suffices to range over
 all-unicyclic graphs; a unicyclic component is a cycle with a rooted tree
 at each cycle vertex, its independence polynomial is the trace of a
 product of $2\times2$ transfer matrices with entries the polynomials of the
 rooted trees with the root in or out, and every quantity is a polynomial
 with nonnegative coefficients in those entries, so only coefficientwise
 minimal objects need to be kept at each stage. The final set of minimal
 polynomials is $\{I_{H_n}\}$ for every $3\le n\le18$. Cross-checks: without
 that pruning, all distinct independence polynomials of all-unicyclic
 graphs were enumerated for $n\le11$ (575 for $n=11$), and the coefficientwise
 minima for $n\le7$ agree with a brute-force run of the delivered program
 over all $n^n$ systems. Missing: a proof for $n\ge10$; the first coefficient
 not covered by Proposition~\ref{fbl:prop:small} is that of $z^4$ at
 $n=10$, where $[z^4]I_{H_{10}}=18$.
\item\label{fbl:fq:critical-block} \textbf{Critical block minimization}
 (Conjecture~\ref{fbl:conj:critical-block}). Proved: the coefficients of
 $z^0,z^1,z^2$ and of every $z^r$ with $r>n/q$, hence the whole conjecture
 for $n\in\{q,2q\}$ (Proposition~\ref{fbl:prop:small}); for $k=1$ it is
 the case $3\mid n$ of item~\ref{fbl:fq:pseudoforest}, so it holds for
 $n\le18$. Missing: every case with $k\ge2$ and $n\ge3q$, beginning with
 the coefficient of $z^3$ at $(k,n)=(2,15)$, which is $125$ for $3K_5$.
\item\label{fbl:fq:sap1} \textbf{Exact strength of $\SAP_1$}
 (Question~\ref{fbl:q:sap1}). Proved: $\SAP_1\Rightarrow\Cm_2$
 (Theorem~\ref{fbl:thm:choice-sandwich}) and, by
 Proposition~\ref{fbl:prop:three}, $\SAP_1\Rightarrow\Cm_3$. So $\SAP_1$ is
 equivalent to choice for pairs only if $\Cm_2$ implies $\Cm_3$ over ZF;
 the write did not consult the literature on implications among the
 principles $\Cm_m$. [Independent check, 5~October 2026: $\Cm_2$ does not
 imply $\Cm_3$ over ZF, so $\SAP_1$ is strictly stronger than choice for
 pairs. Jech~\cite[Theorem~7.16, p.~111]{Jech}: if $m,n$ do not satisfy
 Mostowski's condition~(S), that $n$ has no decomposition $n=p_1+\dots+p_s$
 into primes all larger than~$m$, then there is a model of set theory in
 which $\Cm_k$ holds for every $k\le m$ but $\Cm_n$ fails. For $m=2$, $n=3$
 the condition fails, since $3$ is itself a prime larger than~$2$. The model
 is a permutation model, transferred to ZF as in Jech's Example~7.13, and
 Jech's notes to Chapter~7 credit this independence result to Mostowski
 (1945, for ZFA). In that model $\Cm_2$ holds and $\Cm_3$ fails, so
 $\SAP_1$ fails by Proposition~\ref{fbl:prop:three}. The lead is the
 check's; Jech's statement was read when the check was recorded.] The
 source's sketch (colouring the components of a
 partial functional graph ``appears to require choosing phases or markers,
 but much less than arbitrary propositional compactness'') is a heuristic.
 Missing: any upper bound on $\SAP_1$ below $\BPI$ other than
 well-orderability of the domain (Theorem~\ref{fbl:thm:wellorderable}).
\item\label{fbl:fq:hierarchy} \textbf{The hierarchy $\SAP_k$}
 (Question~\ref{fbl:q:hierarchy}). Proved: $\BPI\Rightarrow\SAP_k\Rightarrow
 \bigwedge_{2\le m\le k+1}\Cm_m\wedge\Cm_3$, and $\BPI$ implies a principle
 $\SAP^{+}_{\mathrm{fin}}$ that implies every $\SAP_k$
 (Proposition~\ref{fbl:prop:finite}). The sentence after
 Theorem~\ref{fbl:thm:choice-sandwich}, ``$\BPI$ is substantially stronger
 than the displayed bounded finite-choice principles'', is given without
 proof or reference; read as ``$\bigwedge_{2\le m\le k+1}\Cm_m$ does not
 imply $\BPI$ over ZF'', it was not checked by the write, and the open
 problem of locating $\SAP_k$ in the gap presupposes it. L\"auchli's
 paper~\cite{Lauchli}, in the source's bibliography but cited nowhere in
 its text, bears on the last clause: by its abstract (seen by the write
 only in a search-engine excerpt of the publisher's page) it assigns to
 every Boolean algebra a graph whose finite subgraphs are $3$-colourable and
 whose $3$-colourability yields a prime ideal. Such graphs need not be
 conflict graphs of $k$-bounded systems, and the write did not examine
 whether the construction can be adapted. Missing: whether some $\SAP_k$
 implies $\SAP_{k+1}$ or $\BPI$, and their places in the Howard--Rubin
 catalogue.
\item\label{fbl:fq:countable} \textbf{The countable-colour theorem}
 (Question~\ref{fbl:q:countable}). The statement is Theorem~4 of de Bruijn
 and Erd\H{o}s (read; note in Section~\ref{fbl:sec:dbe-credit}), whose proof
 applies their Theorem~3 to each class $S_k$ under the Axiom of Choice.
 Re-scoped by Proposition~\ref{fbl:prop:finite}: over ZF it lies between
 $\BPI$ and choice for families of nonempty finite sets. Missing: its exact
 position, and whether it implies any $\SAP_k$.
\item\label{fbl:fq:exact} \textbf{Exact $A_r(n,k)$}
 (Question~\ref{fbl:q:exact}). Known exactly: $r=2$
 (Theorem~\ref{fbl:thm:pairs}), $(r,k)=(3,1)$
 (Theorem~\ref{fbl:thm:k1-triples}), the zero range
 (Corollary~\ref{fbl:cor:threshold}), and, from item~\ref{fbl:fq:pseudoforest},
 every $r$ for $k=1$, $n\le18$ (computation). Missing: everything else, in
 particular the $O(n^{r-2})$ term of \eqref{fbl:eq:asymptotic}.
\item\label{fbl:fq:stability} \textbf{Stability and supersaturation}
 (Question~\ref{fbl:q:stability}). Proved: rigidity at the critical size
 (Theorem~\ref{fbl:thm:critical}). Missing: any stability or
 supersaturation statement.
\item\label{fbl:fq:enumeration} \textbf{Enumeration of minimizers}
 (Question~\ref{fbl:q:enumeration}). Proved: the critical count
 \eqref{fbl:eq:critical-count}. The expectation beyond threshold (``a
 mixture of tournament counts and sparse stability theory'') is a
 heuristic. Missing: any count beyond threshold.
\item\label{fbl:fq:certificate} \textbf{Linear-time certificates for
 extremizers} (the third item of the list in
 Section~\ref{fbl:sec:algorithms}). Proved: linear-time checks for critical
 systems and for minimizers of $A_2(n,k)$ (note there). Missing: a check
 for claimed minimizers of $A_r(n,k)$, $r\ge3$, or of $\alpha(G_F)$ above
 the critical size, that does not search the $r$-subsets.
\item\label{fbl:fq:borel} \textbf{Borel partition number}
 (Question~\ref{fbl:q:borel}). Proved: perfect safe sets for countable
 sections (Theorem~\ref{fbl:thm:perfect}), which is not a partition (the
 source's own caveat). Missing: everything asked.
\item\label{fbl:fq:measurable} \textbf{Measurable positive safe sets}
 (Question~\ref{fbl:q:measurable}). Proved: almost surely safe countable
 samples (Theorem~\ref{fbl:thm:random}). Missing: everything asked.
\item\label{fbl:fq:regularity} \textbf{Regularity strength in ZF}
 (Question~\ref{fbl:q:regularity}). The proofs of
 Theorems~\ref{fbl:thm:perfect} and~\ref{fbl:thm:random} are written in
 ZFC. Missing: the fragments of choice they need, and their separation from
 $\SAP_k$.
\item\label{fbl:fq:cardinal} \textbf{Cardinal blacklists}
 (Question~\ref{fbl:q:cardinal}). Proved here: nothing; the classical
 free-set literature for set mappings of infinite order (note in
 Section~\ref{fbl:sec:dbe-credit}) was not surveyed by the source or the
 write. Missing: everything asked.
\item\label{fbl:fq:random} \textbf{Random blacklist systems}
 (Question~\ref{fbl:q:random}). The worst case is
 $\lceil n/(2k+1)\rceil$ (Theorem~\ref{fbl:thm:alpha}; the question prints
 $n/(2k+1)$). That random sparse systems ``should have substantially larger
 safe sets'' is a heuristic of the source, not proved. Missing: everything
 asked.
\item\label{fbl:fq:hypergraph} \textbf{Hypergraph avoidance}
 (Question~\ref{fbl:q:hypergraph}). Missing: everything asked.
\item\label{fbl:fq:game} \textbf{Game and online forms}
 (Question~\ref{fbl:q:game}). Missing: everything asked.
\item\label{fbl:fq:formal} \textbf{Formalization and sequences.} The plan
 M1--M7 of Section~\ref{fbl:sec:formal} and the OEIS subproject above are
 proposals; nothing is formalized or submitted.
\end{enumerate}

\begin{proposition}[{[write]} Small cases, 5~October 2026]
\label{fbl:prop:small}
\begin{enumerate}[label=(\alph*)]
\item Let $n\ge3$ and let $G$ be an $n$-vertex pseudoforest. Then
 $i_r(G)\ge i_r(H_n)$ for $r\le3$ and for every $r>\deg I_{H_n}$. Here
 the degree of $I_{H_n}$ is $a$ if $n=3a$, and $a+1$ if $n=3a+1$ or
 $n=3a+2$. Consequently Conjecture~\ref{fbl:conj:pseudoforest-poly} holds for
 $3\le n\le9$.
\item Let $k\ge0$, $q=2k+1$, $n=tq$ with $t\ge1$, and let $G$ be the
 conflict graph of a $k$-bounded system on $n$ points. Then
 $i_r(G)\ge i_r(tK_q)$ for $r\le2$ and for every $r\ge t+1$. Consequently
 Conjecture~\ref{fbl:conj:critical-block} holds for $n\in\{q,2q\}$.
\end{enumerate}
\end{proposition}

\begin{proof}
(a) By the first paragraph of the proof of
Theorem~\ref{fbl:thm:k1-triples}, the $n$-vertex pseudoforests are exactly
the conflict graphs of one-bounded systems on $n$ points. For every graph,
$i_0=1$ and $i_1=n$. By Lemma~\ref{fbl:lem:edge-budget} with $k=1$ and
$S=X$, $G$ has at most $n$ edges, so $i_2(G)=\binom n2-e(G)\ge\binom n2-n$;
each component $K_3$, $C_4$, $C_5$ of $H_n$ has as many edges as vertices,
so $H_n$ has exactly $n$ edges and $i_2(H_n)=\binom n2-n$. By
Theorem~\ref{fbl:thm:k1-triples}, $i_3(G)\ge A_3(n,1)$, and the graphs that
the theorem names as attaining $A_3(n,1)$ are $aK_3$, $(a-1)K_3\sqcup C_4$
and $(a-1)K_3\sqcup C_5$, that is, $H_n$; so $i_3(G)\ge i_3(H_n)$. The
independence polynomial of a disjoint union is the product of those of the
components, so $I_{H_n}$ is $(1+3z)^a$, $(1+3z)^{a-1}(1+4z+2z^2)$ or
$(1+3z)^{a-1}(1+5z+5z^2)$, of degree $a$, $a+1$ and $a+1$ (all coefficients
of the factors are positive), and $i_r(H_n)=0\le i_r(G)$ for $r$ beyond
the degree. For $n=3,\dots,9$ the degrees are $1,2,2,2,3,3,3$, at most
$3$, so every $r$ falls under one of the two cases.

(b) For every graph, $i_0=1$ and $i_1=n$. By
Lemma~\ref{fbl:lem:edge-budget}, $e(G)\le kn$, so $i_2(G)\ge\binom n2-kn$;
and $tK_q$ has $t\binom q2=tqk=kn$ edges, so $i_2(tK_q)=\binom n2-kn$. An
independent set of $tK_q$ has at most one vertex in each of the $t$
cliques, so $i_r(tK_q)=0\le i_r(G)$ for $r\ge t+1$. For $t\le2$ every
$r$ satisfies $r\le2$ or $r\ge t+1$.
\end{proof}

\begin{proposition}[{[write]} $\SAP_1$ gives choice for triples, 5~October 2026]
\label{fbl:prop:three}
Over ZF, $\SAP_k$ implies $\Cm_3$ for every $k\ge1$. In particular
$\SAP_1\Rightarrow\Cm_2\wedge\Cm_3$. (For $k\ge2$ this is already part of
Theorem~\ref{fbl:thm:choice-sandwich}; the new case is $k=1$.)
\end{proposition}

\begin{proof}
Let $\cA$ be a family of $3$-element sets. For $A\in\cA$ let $\Sigma(A)$ be
the set of permutations of $A$ of order $3$; it has exactly two elements,
$\sigma$ and $\sigma^{-1}$, and is defined from $A$ without any choice. Put
\[
 X=\{(A,\sigma,a):A\in\cA,\ \sigma\in\Sigma(A),\ a\in A\},\qquad
 F(A,\sigma,a)=\{(A,\sigma,\sigma(a))\}.
\]
Since $\sigma$ has no fixed point, $(A,\sigma,a)\notin F(A,\sigma,a)$ and
$\lvert F(x)\rvert=1\le k$. On each fibre $\{A\}\times\{\sigma\}\times A$
the arcs $a\to\sigma(a)\to\sigma^2(a)\to a$ join every two of its three
points, so every fibre is a clique of the conflict graph. $\SAP_k$ gives
$c:X\to2k+1$ with safe colour classes; $c$ is injective on each fibre. Let
$s(A,\sigma)$ be the point $a\in A$ for which $c(A,\sigma,a)$ is least.
Define $g(A)=s(A,\sigma)$ if $s(A,\sigma)=s(A,\sigma^{-1})$, and otherwise
let $g(A)$ be the unique element of $A\setminus\{s(A,\sigma),
s(A,\sigma^{-1})\}$. The definition is symmetric in $\sigma$ and
$\sigma^{-1}$, so $g(A)$ depends only on $A$ and $c$, and $g(A)\in A$. Thus
$g$ is a choice function on $\cA$, defined in ZF from the single
witness $c$. The second sentence follows with
Theorem~\ref{fbl:thm:choice-sandwich} for $m=2$.
\end{proof}

Let $\SAP_{\mathrm{fin}}$ be the statement that for every set $X$ and every
$F:X\to[X]^{<\omega}$ with $x\notin F(x)$ there is $c:X\to\omega$ whose
colour classes are safe (the conclusion of Theorem~4
of~\cite{deBruijnErdos}), $\SAP^{+}_{\mathrm{fin}}$ the same statement with
the extra requirement $c(x)\le2\lvert F(x)\rvert$ for every $x$, and
$\Cm_{\mathrm{fin}}$ the axiom of choice for families of nonempty finite
sets.

\begin{proposition}[{[write]} Countable colours over ZF, 5~October 2026]
\label{fbl:prop:finite}
Over ZF,
\[
 \BPI\Longrightarrow\SAP^{+}_{\mathrm{fin}}\Longrightarrow\SAP_{\mathrm{fin}}
 \Longrightarrow\Cm_{\mathrm{fin}},
 \qquad\text{and}\qquad
 \SAP^{+}_{\mathrm{fin}}\Longrightarrow\SAP_k\quad\text{for every }k\ge1.
\]
\end{proposition}

\begin{proof}
\emph{$\BPI\Rightarrow\SAP^{+}_{\mathrm{fin}}$.} Let $F:X\to[X]^{<\omega}$,
$b(x)=\lvert F(x)\rvert$ and $L(x)=\{0,1,\dots,2b(x)\}$. Take propositional
variables $p_{x,j}$ for $x\in X$ and $j\in L(x)$, and the clauses
$\bigvee_{j\in L(x)}p_{x,j}$ (a finite disjunction),
$\neg(p_{x,i}\wedge p_{x,j})$ for $i\ne j$ in $L(x)$, and
$\neg(p_{x,j}\wedge p_{y,j})$ for every conflict edge $\{x,y\}$ and
$j\in L(x)\cap L(y)$. A finite set of these clauses mentions a finite set
$S$ of points. On $S$ consider the system $F_S(x)=F(x)\cap S$; its conflict
graph is the induced graph $G_F[S]$, and $\lvert L(x)\rvert=2b(x)+1\ge
2\lvert F_S(x)\rvert+1$. Theorem~\ref{fbl:thm:list} gives a proper list
colouring of $G_F[S]$ from the lists $L(x)$, and setting $p_{x,j}$ true
exactly when $x$ has colour $j$ satisfies the finite set. $\BPI$ is
equivalent over ZF to the compactness theorem of propositional logic (the
fact used in Theorem~\ref{fbl:thm:dbe}), so all clauses are satisfied by one
valuation; let $c(x)$ be the unique $j\in L(x)$ with $p_{x,j}$ true. Then
$c(x)\le2b(x)$, and two conflicting points get different colours, so every
colour class is safe.

\emph{$\SAP^{+}_{\mathrm{fin}}\Rightarrow\SAP_{\mathrm{fin}}$} is trivial,
and \emph{$\SAP^{+}_{\mathrm{fin}}\Rightarrow\SAP_k$}: for a $k$-bounded
system, $c(x)\le2k$, so $c$ maps into $2k+1$.

\emph{$\SAP_{\mathrm{fin}}\Rightarrow\Cm_{\mathrm{fin}}$.} This is the
clique-fibre argument of Theorem~\ref{fbl:thm:choice-sandwich} without the
size bound. Let $\cA$ be a family of nonempty finite sets, $X=\{(A,a):
A\in\cA,\ a\in A\}$ and $F(A,a)=\{(A,b):b\in A\setminus\{a\}\}$, a finite set
not containing $(A,a)$. Each fibre $\{A\}\times A$ is a clique of the
conflict graph, so a colouring $c:X\to\omega$ with safe classes is
injective on it, and choosing from $A$ the element of least colour defines
a choice function on $\cA$ in ZF.
\end{proof}

\begin{writenote}[Independent check of the write, 5 October 2026]
An adversarial check made by the intake after the write (commit
\file{bf70d7db4}) found no item needing a correction. It reread the proofs
of Propositions~\ref{fbl:prop:small}, \ref{fbl:prop:three}
and~\ref{fbl:prop:finite}, every implication of the last included, and the
note completing the minimizer step of Theorem~\ref{fbl:thm:k1-triples}; it
confirmed that for $k\ge2$ choice for triples is already the case $m=3$ of
Theorem~\ref{fbl:thm:choice-sandwich}, so Proposition~\ref{fbl:prop:three}
is new exactly for $k=1$. With its own program, by a method different from
the write's transfer matrices (all nonisomorphic trees, and every tree plus
one edge, as the connected pseudoforests; independence polynomials by
bitmask recursion; coefficientwise minima for each order and over all
partitions of $n$; not shipped), it found the coefficientwise minimal
independence polynomial over all $n$-vertex pseudoforests, trees included,
to be exactly $I_{H_n}$ for every $3\le n\le15$. Brute force over all $n^n$
one-bounded systems for $n\le6$, and over all $11^5$ two-bounded systems on
five points (unique minimum $1+5z$), agreed, as did the closed forms
$A_2(n,1)$ and $A_3(n,1)$, the printed table, $[z^4]I_{H_{10}}=18$ and
$[z^3]I_{3K_5}=125$. It checked the account of de Bruijn and Erd\H{o}s
against the same scan (Theorems~1, 3 and~4 and their proofs; three pages,
numbered [371], 372 and~373, with Theorem~4 on p.~372); a web search finds
the variant 369--373 attributed to \emph{Indag.\ Math.}\ in secondary
sources, which neither the check nor the write resolved. Its one lead, that
choice for pairs does not imply choice for triples over ZF, is confirmed
from Jech's book in item~\ref{fbl:fq:sap1}.
\end{writenote}

\section{Formalization plan for ProveIt}\label{fbl:sec:formal}

A compact Lean development can be organized into the following modules.

\begin{enumerate}[label=\texttt{M\arabic*.}]
\item \textbf{Blacklist and conflict graph.}  Finite types, the directed
relation $y\in F(x)$, symmetrization, and equivalence between safe sets and
independent sets.
\item \textbf{Edge budget.}  Map each conflict edge to a witnessing arc,
prove $e(S)\le\sum_{x\in S}|F(x)\cap S|$, and derive a low-degree vertex.
\item \textbf{Peeling and list colouring.}  Formalize the variable-budget
elimination order and greedy reverse colouring.
\item \textbf{Sharp constructions.}  Cyclic tournaments, near-regular even
orientations, block sums, and the exact independence-number theorem.
\item \textbf{Critical equality.}  Import or prove the finite Tur\'an equality
case, then derive the regular-tournament classification and counting formula.
\item \textbf{Spectrum.}  Double counting for pairs, Bonferroni bounds, and
the pseudoforest triple formula.
\item \textbf{Foundations interface.}  State the infinite colouring as an
axiom parameterized by a compactness principle; keep the finite theorem free
of classical-choice dependencies.
\end{enumerate}

The finite core should require only ordinary finite combinatorics from
Mathlib.  The descriptive-set-theoretic perfect-set theorem is a separate,
substantially larger formalization project.

\section{Assumption, novelty, and claim ledger}\label{fbl:sec:ledger}

\begin{longtable}{p{0.26\textwidth}p{0.18\textwidth}p{0.47\textwidth}}
\toprule
Claim or ingredient & Status & Notes\\ \midrule
$2k+1$-colouring of arbitrary infinite systems & Classical &
De Bruijn--Erd\H{o}s (1951), via compactness; their paper adopts AC.\\
Finite list theorem and exact safe number & Proved here &
Elementary edge-budget and greedy arguments; no priority claim after an
exhaustive search.\\
Critical regular-tournament classification & Proved here &
Uses Tur\'an's equality theorem and saturation of the arc budget.\\
Exact critical enumeration & Proved here &
Immediate from the classification and labelled regular tournament count.\\
Exact pair spectrum & Proved here &
Edge count plus an explicit circulant construction.\\
Fixed-$r$ two-term asymptotic & Proved here &
Union bound and second Bonferroni inequality.\\
Exact $A_3(n,1)$ & Proved here &
Pseudoforest reduction and a component cost inequality.\\
Choice sandwich & Proved here &
$\BPI$ upper bound by compactness; bounded finite-choice lower bounds by
clique fibres.\\
Perfect safe set & Classical tools, derived here &
Kuratowski--Ulam plus Mycielski; the application to blacklist symmetry is
spelled out.\\
Random countable safe set & Proved here from standard measure theory &
Tonelli plus a countable union of null events.\\
Coefficientwise minimization conjectures & Open &
Supported only by the proved coefficients and small exhaustive checks.\\
Exact choice strength of $\SAP_k$ & Open &
The central foundational question proposed by the report.\\
\bottomrule
\end{longtable}

\section{Conclusion}

Sparse symmetry-avoidance is deliberately simple to state:
\begin{quote}
Each point names at most $k$ points it refuses to pair with.  How large a
mutually accepting set must exist, and how much choice is needed to organize
all points into finitely many such sets?
\end{quote}
The finite answer begins with a complete theorem: the exact denominator is
$2k+1$, the safe-$r$ threshold is $(r-1)(2k+1)+1$, and the final obstructions
are precisely blocks of regular tournaments.  The quantitative theory then
opens into safe-set spectra, pseudoforest formulas, stability, and
enumeration.

The infinite answer is more revealing.  Local finite colourability alone
does not specify which compactness principle may be used to assemble a
global colouring.  The principle $\SAP_k$ is strong enough to imply genuine
finite-choice axioms, weak enough that its exact position below $\BPI$ is
unclear, and automatically true on well-orderable sets.  Meanwhile Borel and
measurable regularity bypass the partition question and produce perfect or
random countably infinite safe sets.  This combination of exact finite
combinatorics, weak choice, Freiling symmetry, and definability makes the
problem a natural addition to the ProveIt research program and a plausible
point of contact with Glazer's published interests.

\appendix

\section{A self-contained compactness lemma for well-orderable languages}

For completeness, we isolate the logical fact used in
\cref{fbl:thm:wellorderable}.

\begin{lemma}\label{fbl:lem:wellordered-compactness}
In ZF, let $V$ be a well-orderable set of propositional variables and let
$T$ be a set of formulas over $V$.  If every finite subset of $T$ is
satisfiable, then $T$ is satisfiable.
\end{lemma}

\begin{proof}
The set of finite formulas over a well-orderable alphabet is well-orderable.
Enumerate it as $\angles{\varphi_\alpha:\alpha<\theta}$.  Starting with
$T_0=T$, recursively let
\[
 T_{\alpha+1}=
 \begin{cases}
 T_\alpha\cup\{\varphi_\alpha\},&
    \text{if this is finitely satisfiable},\\
 T_\alpha\cup\{\neg\varphi_\alpha\},&\text{otherwise},
 \end{cases}
\]
and take unions at limit stages.  At least one successor choice preserves
finite satisfiability: if both failed, finite witnesses from the two failures
would combine into a finite unsatisfiable subset of $T_\alpha$.  The union
$T^*$ is complete and finitely satisfiable.  Set a variable $p$ to true
exactly when $p\in T^*$.  Induction on formula complexity shows that every
member of $T^*$ is true under this valuation, and hence every member of $T$
is true.
\end{proof}

\section{Reproduction notes}

The archive contains:
\begin{itemize}
\item \texttt{article.tex}: the complete source of this report;
\item \texttt{article.pdf}: the compiled report;
\item \texttt{verification.py}: the finite exhaustive checks described in
\cref{fbl:sec:questions} and the computational section;
\item \texttt{README.txt}: a short build and verification guide.
\end{itemize}

To run the finite checks:
\begin{verbatim}
python3 verification.py
\end{verbatim}
The script uses only the Python standard library.  The $n=7$ exhaustive
check is the largest one and may take several seconds to a few minutes,
depending on hardware.

\begin{writenote}[Names in this appendix, 5 October 2026]
The list above describes the delivered archive. There the source and the
PDF were named \file{From_Finite_Blacklists_to_Freiling_Symmetry.tex} and
\file{.pdf}, not \file{article.tex} and \file{article.pdf}; the archive
also held \file{verification_output.txt} and \file{SHA256SUMS.txt}. In this
report the source is \file{article.tex}, the program is
\file{code/verification.py}, its recorded output is
\file{data/verification_output.txt}, and the README replaces
\file{README.txt}. The checks are described in
Section~\ref{fbl:sec:computation}, not in \cref{fbl:sec:questions}. The
program writes no file and can be run from any directory; on Windows use
\texttt{py} for \texttt{python3}. The whole run took about 6~seconds at
the write.
\end{writenote}

\begin{thebibliography}{99}

\bibitem{GlazerThesis}
E.~Glazer,
\emph{Interactions Among Principles of Choice, Forcing, and Measure Theory},
Ph.D. dissertation, Harvard University, 2023.
\url{https://dash.harvard.edu/entities/publication/12fe6154-3d3a-47ec-8576-afc35fae7e36}.

\bibitem{GlazerBox}
E.~Glazer,
``A choiceless box game paradox,''
\emph{arXiv:2211.10474}, 2022.
\url{https://arxiv.org/abs/2211.10474}.

\bibitem{GlazerTennenbaum}
E.~Glazer,
``A Topological Tennenbaum Theorem,''
\emph{arXiv:2311.13699}, 2023.
\url{https://arxiv.org/abs/2311.13699}.

\bibitem{Freiling}
C.~Freiling,
``Axioms of symmetry: throwing darts at the real number line,''
\emph{Journal of Symbolic Logic} \textbf{51} (1986), no.~1, 190--200.
\doi{10.2307/2273955}.

\bibitem{deBruijnErdos}
N.~G. de Bruijn and P.~Erd\H{o}s,
``A colour problem for infinite graphs and a problem in the theory of
relations,''
\emph{Proceedings of the Koninklijke Nederlandse Akademie van Wetenschappen,
Series A} \textbf{54} (1951), 371--373; also
\emph{Indagationes Mathematicae} \textbf{13} (1951), 369--373.
\url{https://pure.tue.nl/ws/files/4237754/597497.pdf}.

\bibitem{Turan}
P.~Tur\'an,
``On an extremal problem in graph theory,''
\emph{Matematikai \ és Fizikai Lapok} \textbf{48} (1941), 436--452.

\bibitem{Cowen1990}
R.~Cowen,
``Two hypergraph theorems equivalent to the Boolean prime ideal theorem,''
\emph{Notre Dame Journal of Formal Logic} \textbf{31} (1990), no.~2,
232--240.
\doi{10.1305/ndjfl/1093635418}.

\bibitem{Cowen1998}
R.~Cowen,
``A compactness theorem for infinite constraint satisfaction,''
\emph{Reports on Mathematical Logic} \textbf{32} (1998), 97--107.

\bibitem{Lauchli}
H.~L\"auchli,
``Coloring infinite graphs and the Boolean prime ideal theorem,''
\emph{Israel Journal of Mathematics} \textbf{9} (1971), 422--429.

\bibitem{HowardRubin}
P.~Howard and J.~E. Rubin,
\emph{Consequences of the Axiom of Choice},
Mathematical Surveys and Monographs 59, American Mathematical Society, 1998.

\bibitem{Jech}
T.~Jech,
\emph{The Axiom of Choice},
North-Holland, 1973.

\bibitem{Mycielski}
J.~Mycielski,
``Independent sets in topological algebras,''
\emph{Fundamenta Mathematicae} \textbf{55} (1964), no.~2, 139--147.
\doi{10.4064/fm-55-2-139-147}.

\bibitem{Kechris}
A.~S. Kechris,
\emph{Classical Descriptive Set Theory},
Graduate Texts in Mathematics 156, Springer, 1995.

\bibitem{Bollobas}
B.~Bollob\'as,
\emph{Modern Graph Theory},
Graduate Texts in Mathematics 184, Springer, 1998.

\bibitem{ProveIt}
V.~Reshetnikov,
\emph{ProveIt}, public research repository, accessed 5 October 2026,
commit \texttt{a763feee1f3703dc710d9b810b1de66168f1e266}.
\url{https://github.com/VladimirReshetnikov/ProveIt}.

\bibitem{ProveItBoxes}
V.~Reshetnikov,
``Countable Information, Uncountably Many Boxes,''
repository report in \texttt{ProveIt/SetTheory/Cardinals/docs/reports/
ordinals-and-order-types/measurable-box-games}, 2026.
\url{https://github.com/VladimirReshetnikov/ProveIt/tree/main/SetTheory/Cardinals/docs/reports/ordinals-and-order-types/measurable-box-games}.

% Entries added by the ProveIt write phase (batch 114, 5 October 2026).
\bibitem{Erdos1950}
P.~Erd\H{o}s,
``Some remarks on set theory,''
\emph{Proceedings of the American Mathematical Society} \textbf{1} (1950),
127--141. [Added by the write, as cited by~\cite{deBruijnErdos}; not read.]

\bibitem{Sierpinski}
W.~Sierpi\'nski,
\emph{Hypoth\`ese du continu}, 2nd ed., Chelsea, New York, 1956.
[Added by the write, as cited by~\cite{Freiling}; not read.]

\bibitem{Hajnal}
A.~Hajnal,
``Proof of a conjecture of S.~Ruziewicz,''
\emph{Fundamenta Mathematicae} \textbf{50} (1961), 123--128.
[Added by the write; data from a search listing; not read.]

\bibitem{OEIS}
OEIS Foundation Inc.,
\emph{The On-Line Encyclopedia of Integer Sequences}, entries A133686 and
A007079, \url{https://oeis.org}, accessed 5 October 2026. [Added by the
write.]

\end{thebibliography}

\end{document}
